\begin{abstract}
We study a one-parameter family of partition functions $\{Z_m(y)\}_{m\ge 0}$ governed by a quartic spectral curve.
We prove that the ordinary generating function $Z(t,y)=\sum_{m\ge 0} Z_m(y)t^m$ is a reduced rational function whose denominator has minimal degree~$4$, implying a minimal linear recurrence of order~$4$ and a fixed Hankel rank~$4$ for all parameters $y$ outside a finite degeneracy set.
We then compute and completely factor the discriminant of the quartic spectral polynomial $\Pi(\lambda,y)$, obtaining
\[
\Disc_\lambda\Pi(\lambda,y)=-y(y-1)(256y^3+411y^2+165y+32),
\]
which determines the branch locus exactly.
The cubic factor has a unique real root $\yLY\in(-1.13446,-1.13445)$, which is a nontrivial branch point where two subdominant spectral roots collide and transition from real to complex conjugate; the full equimodular zero-accumulation geometry, governed by the Beraha--Kahane--Weiss theorem, is explicitly controlled by this branch locus.
Introducing an exponential family via $Z_m(e^s)/Z_m(1)$, we prove that the limiting log-moment generating function is controlled by the dominant spectral branch $\lambda_+(y)$, and we derive closed-form constants for a law of large numbers and a central limit theorem: mean $\mu=5/18$ and variance $\sigma^2=67/972$.
Finally, we show that the quartic spectral curve is birationally equivalent to the conductor-$37$ elliptic curve $\xi^2+\xi=\lambda^3-\lambda$ (Cremona label~$37\mathrm{a}1$), and that the branch locus arises as the ramification image of a degree-$4$ morphism on this curve.
\end{abstract}

\begin{sloppypar}
\textbf{Keywords:} partition-function zeros; discriminant factorization; spectral curve; Hankel rank; free energy; central limit theorem; linear recurrence; elliptic curve; Beraha--Kahane--Weiss theorem
\end{sloppypar}

\tableofcontents
\newpage

